\section{Implementation and independent certificate replay}
\label{sec:implementation}

\subsection{A small, opt-in native change}

The patch changes four production modules and adds two test modules under
\path{Topology/UnknotRecognition/fast}.  Existing defaults remain
canonical.  The optional argument
\begin{center}
\code{forest\_strategy='deferred'}
\end{center}
is accepted by the core producer, finite bundle-boundary producer, and
whole-cut wrapper.  Unknown strategy values raise \code{ValueError}.
The argument is propagated through the supported core, augmented-core,
and finite-kernel fallback calls.

In the core producer, deferred discovery builds the full occurrence
matrix, runs the existing column queue, and chooses a spanning forest
avoiding the recorded pivot columns.  It checks that the original
component partition is unchanged.  If peeling stops, the result remains
\code{INCONCLUSIVE}, with reason
\code{all\_primal\_forests\_raw\_obstruction} and the remaining two-handle
identifiers.  On success it emits the existing
\code{normal-core-handlebody-v2} schema.  The independent core checker
already accepts a general supplied forest, so its code is unchanged.

The logical producer is short:
\begin{Code}
rows = source_occurrences_without_a_primal_forest()
trace, residual = column_leaf_peeling(rows)
if residual:
    return INCONCLUSIVE, restricted_stopping_witness(residual)
pivots = source_graph_columns_used_by(trace)
forest = spanning_forest(source_graph_without(pivots))
assert original_components_are_preserved(forest)
return certificate(forest_first=True, forest=forest, trace=trace)
\end{Code}
This pseudocode describes one supported finite matrix.  The full wrapper
still has conservative branches, port conditions, and bundle lifting
requirements outside this local argument.

\subsection{Why the whole-boundary change needs a new schema}

Core success is insufficient when a whole cut is reconstructed by gluing
parallelity bundles.  The attaching annuli add rows and may destroy
leaves that were available in the isolated core.  Deferred boundary
discovery therefore first includes every supported annular row, using
the authenticated closed core paths and dedicated auxiliary base
columns.  It then peels this full finite matrix and chooses the source
forest avoiding its source pivots.

The old boundary schemas reconstruct a canonical forest.  Reinterpreting
them as permitting an arbitrary supplied forest would change an existing
certificate's meaning.  The patch instead introduces
\code{bundle-boundary-kernel-v3}; earlier schemas retain their exact
validation policy.  The outer whole-cut certificate schemas need not
change because their kernel field already dispatches by kernel schema.

The v3 checker independently decodes and validates every supplied forest
edge.  It rejects booleans masquerading as integer identifiers, malformed
hexadecimal strings, negative or out-of-range handles, duplicates, loops,
cycles, and a failure to span the original graph components.  It uses
authenticated endpoints to reconstruct the component partition, rather
than calling the producer's forest routine.  It then reconstructs raw
post-forest source and annular occurrences, validates the cancellation
trace, and recomputes the component census, base families, and bridge
information.

The component identifiers remain the minimum source vertex labels, so
choosing another forest does not silently rename component ownership.
The tree is removed from the surviving-column census after discovery;
forgetting that removal would preserve a valid cancellation trace but
give an incorrect genus.  Both steps are covered by direct kernel and
whole-cut replay tests.

\subsection{Checking independence and hostile certificates}

The new tests monkeypatch producer entry points, the producer's forest
helper, the peeling routine, and a census helper to raise if called while
verifying a saved positive packet.  Verification still succeeds.  This
tests the intended separation at the modified interfaces; it is not a
claim that source preparation and all geometry code are formally
independent implementations.

Mutation cases remove a tree edge, duplicate an edge, insert invalid
labels, remove a required cancellation, select a forest edge as a
cancellation pivot, or relabel a v3 packet as v2.  All are rejected.
Further tests exercise callbacks, exhausted operation limits, source
relabelings, an exact literal oracle comparison, and a normal-coordinate
scale factor $2^{128}$.  The scaling test checks that binary multiplicity
does not force literal normal-sheet expansion in the changed path; it
does not establish an asymptotic runtime bound by measurement alone.

\subsection{The standalone exceptional solver}

The research operator \code{exceptional\_forest.py} depends only on
Python's standard library.  It uses the attributed, unchanged
\code{degree\_two\_backend.py} from the incoming dual-normalization
package for validation, strongly connected components, classical matroid
intersection, forest extension, and final peeling.  Its new code computes
the exceptional interface and signatures, performs Horn closure,
enumerates partial partitions, and clones overlapping candidate
memberships.

The input has one optional dual-edge record per matrix row.  A null
record is ineligible for the selected tree, and loops cannot be selected.
The root is vertex zero.  Sparse row entries retain exact positive integer
counts on input, although the solver caps them at two for singleton
decisions.  The three statuses mean:
\begin{description}[style=nextline]
\item[\code{SUCCESS}] A final spanning tree and complete peel sequence
have been constructed.  A separate finite checker can verify both.
\item[\code{NO\_FOREST}] The complete finite search has excluded every
tree in the specified graph/matrix problem.  This is not a topological
negative certificate.
\item[\code{INCONCLUSIVE}] A work/partition budget or implementation
limit interrupted the search.  Exhaustion of a resource is never treated
as exhaustive mathematical failure.
\end{description}

No native extraction adapter for this operator is included.  Such an
adapter should use the existing authenticated dual endpoints and full
raw source occurrences, feed any successful tree into the legal top
reduction, and obtain ordinary native replay before claiming a new
topological success.  The finite solver's tests deliberately do not
assume geometric realizability of random matrices.

\subsection{Integration policy}

The direct patch preserves canonical behavior by default.  This is a
practical decision: the stronger deferred test adds work on many inputs
that the canonical policy already settles.  A later adaptive policy could
try canonical peeling first and defer only when needed, but should reuse
authenticated rows and avoid rebuilding expensive geometry.  We do not
include that additional optimization without measuring its full cost.

The artifact contains both a patch against the pinned commit and an
overlay of the six changed files.  Apply one representation, not both.
It also contains source hashes, the original baseline versions of the
four production modules for comparison, exact new packets, and the audit
drivers.  Existing source validators, the normal-orbit construction,
dual breadth-first policy, and the outer knot verdict logic retain their
own contracts; local completeness must not be read as a claim about
every branch of the whole wrapper.
